\paragraph{Theorem E.2.16 (Proved: the full complement module and its peripheral map).}
Let $K=\pi_1(E)$ have the marking (E.S1), and put
\[
 R=\mathbb ZK,\qquad Q=K/\langle h\rangle=C_3*C_4,
 \qquad \overline R=\mathbb ZQ.
\]
Use right deck modules; conversion to left modules applies the
involution $g\mapsto g^{-1}$ to every coefficient. Then
\[
 \pi_2(E)\simeq R\oplus\overline R^{\,2}\oplus L,
 \qquad
 L=\langle x,t\mid t(a-1)=t(b-1)=0,
                         \ x(1+b)+3t=0\rangle_{\overline R}.
 \tag{E.C1}
\]
Equivalently, replacing $x$ by $x+t$ and eliminating $t$ gives
\[
 L\simeq\overline R\big/
       \big((b^2-1)\overline R+(1+b)(a-1)\overline R\big).
 \tag{E.C2}
\]
There is a nonsplit exact sequence
\[
 0\longrightarrow\mathbb Z\langle t\rangle
  \longrightarrow L\longrightarrow
  \mathbb Z_{-}\otimes_{\mathbb ZC_4}\overline R
  \longrightarrow0,
 \tag{E.C3}
\]
where $b$ acts as $-1$ on $\mathbb Z_{-}$ and the first term
has trivial $Q$ action. Its extension class is the nonzero
element of an order-two extension group. In the connected
lift of the peripheral $T^5$ to the universal cover of $E$,
the degree-two classes $uw,ud,wd$ map respectively to
\[
                           t,\quad0,\quad0.
 \tag{E.C4}
\]
Thus the central torsion summand of the module is annihilated
by $h-1$ itself. Formula (E.C1) is a full integral module
calculation, including the central value $h=1$.  The complete
splice extension is determined separately in E.2.17, and the
smooth type by full framed handles in E.2.21; neither conclusion
is inferred from this module calculation alone.

There is also a specified chain-level peripheral comparison.
For the splitting constructed below, let $e$ generate the
free $R$ summand and let $y,z$ generate the two copies of
$\overline R$. The projection of the peripheral degree-two
chain map to (E.C1) is
\[
 \begin{array}{c|rrrrrr}
 \text{cell}&gu&gw&gd&uw&ud&wd\\ \hline
 \text{value}&0&0&6t&t&0&0
 \end{array}
 \tag{E.C19a}
\]
and, writing $\Theta$ for the outer base circle with
group image $\mu=ab$, use the cone convention: the
notation $v\Theta$ denotes the suspended cell
$\Theta\times v$, with the interval factor first.
In this convention,
\[
 \begin{array}{c|l}
  g\Theta&eh-2x-6yb+z(5-2b)\\
  u\Theta&x-yb+zb\\
  w\Theta&-y(1+b)+z\\
  d\Theta&x+2t+z(b-1).
 \end{array}
 \tag{E.C19b}
\]
These are values on chains, not assertions that all
the displayed cells are cycles in the cover. In
particular the free coefficient $h$ on $g\Theta$
is nonzero even though the free summand receives
zero from the lifted peripheral homology in (E.C4).
The original cone and its full peripheral map, not
only the map on homology, are part of the certificate.

\begin{proof}
We give the local chain comparison before the global
calculation, because the ordinary homology maps alone would
not determine (E.C1).

\emph{The cusp boundary and its derived attachment.}
Write $A=\mathbb Z[h,h^{-1}]$ and $q=h-1$. Let
$\widehat W$ be the gamma cover used in E.2.12, and let
$\widehat M$ be the corresponding boundary cover. Ordinary
chains with the deck action give finite free $A$ chain
models. The homology of $\widehat W$ is
\[
 H_0=H_1=A/(q),\quad H_2=A\oplus(A/(q))^2,
 \quad H_3=0,\quad H_4=A.
 \tag{E.C5}
\]
For the boundary, start with the torus chain complex
$C= A\otimes\bigwedge\langle g,u,w,d\rangle$,
whose differential contracts $g$ and multiplies by $q$.
The cellular monodromy has degree-one matrix
\[
 F(M_0)g=g-hd,\quad F(M_0)u=u+w,
 \quad F(M_0)w=w,\quad F(M_0)d=d.
 \tag{E.C6}
\]
Its higher matrices are exterior powers. The mapping-torus
complex is the cone of $I-F(M_0)$, with differential
\[
 d(c,e)=\big(dc+(I-F(M_0))e,-de\big).
 \tag{E.C7}
\]
For example $d(g\theta)=hd-q\theta$ and
$d(u\theta)=-w$. Cancelling these two unit pairs leaves
the degree-one classes $u$ and $\theta$, with resolutions
$q$ and a unit times $q^2$. In degree two the four surviving
classes are $uw,ud,w\theta,d\theta$, each with resolution
$q$. In particular
\[
 H_0(\widehat M)=A/(q),\quad
 H_1(\widehat M)=A/(q^2)\oplus A/(q),\quad
 H_2(\widehat M)=(A/(q))^4.
 \tag{E.C8}
\]
Geometrically this is the Heisenberg three-manifold times
a circle: the deck transformation shears the meridian
by the $d$ circle. The homology map to (E.C5) is the
identity in degree zero, projection to $u$ in degree one,
and the primitive projection
$(uw,ud,w\theta,d\theta)\mapsto(U,D,0,0)$ in degree two.
The last assertion uses the actual specialization
(B.5.16) and the literal collapsed meridian tori
(B.5.17), rather than ranks alone.

The missing extension is fixed by a relative intersection
calculation. Let $S_i$ denote the normalized hexagon
component with gamma index $i$. In its standard Picard
basis $H,E_1,E_2,E_3$, the six boundary curves are
\[
 E_1,\ H-E_1-E_2,\ E_2,\ H-E_2-E_3,
                       \ E_3,\ H-E_1-E_3.
 \tag{E.C9}
\]
The first opposite pair is self-glued in the surviving
$u$ direction; the other two pairs join $S_i$ to its
two gamma neighbours. Put $C_0=-H+E_1$, the primitive
generator of the free kernel of the normalization matrix
$B(h)$ in E.2.12. The divisor $S_i$ restricts on its
own normalization to $2C_0$: the total central divisor
is principal, and the two self-gluing preimages cancel
their terms in the normal restriction. It restricts
to $-C_0$ on each neighbouring component and to zero
on every other component. Thus the free part of the
equivariant intersection map is
\[
 H_4(\widehat W)=A\longrightarrow
 H_4(\widehat W,\widehat M)\longrightarrow A,
 \qquad 2-h-h^{-1}=-h^{-1}q^2.
 \tag{E.C10}
\]
Poincare--Lefschetz duality over $A$, with its usual
involution, gives the same coefficient, up to a unit,
for the free part of $H_2(\widehat W)\to
H_2(\widehat W,\widehat M)$. The latter relative group
is $A$: the universal-coefficient sequence in degree
four has $\operatorname{Hom}_A(H_4,A)=A$, zero $H_3$,
and no higher extension terms. Indeed every module
in (E.C5) has a projective resolution of length at
most one. The relative sequence therefore includes
\[
 0\longrightarrow A\xrightarrow{\,\text{unit}\ q^2\,}A
  \longrightarrow A/(q^2)\longrightarrow0.
 \tag{E.C11}
\]
Consequently the boundary map has a unit component in
$\operatorname{Ext}^1_A(A/(q^2),A)=A/(q^2)$.
Every unit of $A/(q^2)=\mathbb Z[\varepsilon]/(\varepsilon^2)$
lifts to a unit $\pm h^n$ of $A$.

Here is why this determines the required low chain map
without discarding a peripheral invariant. All the
modules in (E.C5) and (E.C8) have projective dimension
at most one. Their complexes split in the derived
category into their two-term projective resolutions:
the obstructions to splitting consecutive truncations
are extension groups in degree at least two and vanish.
In degrees relevant to $H_2$ of a mapping cone, a map
between these split complexes consists of the displayed
homology maps and extension components from $H_i$ to
$H_{i+1}$. Normalize the component in (E.C11) by a unit
on the free $H_2$ summand of the target. Its components
to the two torsion $H_2$ summands can then be removed
by target maps $A\to(A/(q))^2$: reduction of the unit
modulo $q$ is $\pm1$. The remaining components from
the $u$ summand are removed by target degree-one
extension maps from $H_1(\widehat W)$ to
$H_2(\widehat W)$. Finally the degree-zero component
is removed using the isomorphism on $H_0$ and a target
extension map from $H_0$ to $H_1$.

Each change just described is a unit on the diagonal
of the split complex and is triangular with extension
entries off the diagonal. It is therefore invertible.
Replacing the split resolutions by the original finite
free complexes represents it by a chain homotopy
equivalence. No realization by a diffeomorphism or a
map of the cusp space is claimed or needed. The change
is applied to the cusp target and its entire incoming
boundary map simultaneously. The other vertices and
all maps into the regular core are unchanged. It thus
gives a chain homotopy equivalence of the global
mapping-cone complex that is the identity on the
external peripheral complex. In particular the
calculation of (E.C4) retains the actual boundary
map. Tensoring with $R$ preserves these comparisons,
since $R$ is free over $A$.

One resulting cusp model in degrees zero through three
has ranks $(1,2,4,2)$, generators
\[
 (1);\quad(g,u);\quad(gu,E,U,D);\quad(R_U,R_D),
\]
and differentials $dg=q$, $d(gu)=qu$,
$dR_U=qU$, $dR_D=qD$. The map (E.C7) to this model
is, through degree two,
\[
 \begin{array}{c|ccccc}
  &g&u&w&d&\theta\\ \hline
  i_1&g&u&0&0&0
 \end{array},\qquad
 \begin{array}{c|rrrrrr}
  &gu&gw&gd&uw&ud&wd\\ \hline
  i_2&gu&-hD&E&U&D&0
 \end{array},
 \qquad i_2(v\theta)=0.
 \tag{E.C12}
\]
The signs and the factor $h$ satisfy the chain equation
directly. At $h=1$, choose
$E=-\overline C_1-\overline C_2$,
$U=\overline C_2+\overline C_3$, and
$D=\overline C_2$; this is precisely (B.5.16).

\emph{Finite caps and the global complex.}
For any marked monodromy
$A=\left(\begin{smallmatrix}1&0\\t&D\end{smallmatrix}\right)$
use the chain lift
$F(A)=\left(\begin{smallmatrix}1&0\\ht&D\end{smallmatrix}\right)$
and its exterior powers. These lifts multiply exactly
and commute with contraction by $qg^*$.

This exterior lift represents the actual equivariant
torus map. The gamma cover is
$\mathbb R_g\times T^3$, and the marked linear map is
$(g,v)\mapsto(g,Dv+gt)$. At every integer value of $g$
the translation is integral, so the vertex torus map
is exactly $D$. In a single unit strip, concentrate
the loop of translations $v\mapsto v+st$ at its
terminal endpoint by an equivariant cellular
approximation relative to the two vertex tori.
The translation prism on the minimal cellular complex
of $T^3$ is exterior multiplication by $t$: its value
on each product torus is the homology class of the
literal translation sweep, and all differentials
in that minimal torus complex are zero. The terminal
vertex has deck coefficient $h$. Consequently
\[
 v\longmapsto\bigwedge D(v),\qquad
 g\wedge v\longmapsto
 g\wedge\bigwedge D(v)+ht\wedge\bigwedge D(v),
 \tag{E.C20}
\]
which is exactly the exterior lift above. There is
only this one strip parameter. Composing two such
actual maps adds the translation loops by
$t_A+D_At_B$, the same strict multiplication law
as $F(A)F(B)=F(AB)$. Thus no uncomputed higher
coherence is being substituted for the torus map.
A regular
boundary with deck image $s$ is the cone of $sI-F(A)$.
This convention gives the group relation
$s\lambda s^{-1}=A\lambda$.

At finite cap $j$, take its adapted first fibre vector
$n_jv_j$, where
\[
 (m_1,n_1,v_1)=(3,1,(1,2,-4,0)),\qquad
 (m_2,n_2,v_2)=(4,-1,(-1,-3,3,0)).
 \tag{E.C13}
\]
It has gamma coordinate one and is fixed by $A_j$.
The other three coordinates form $L_0=\langle u,w,d\rangle$,
with monodromy $D_j=A_j|_{L_0}$. The cap is the
mapping torus of this three-torus, and its complex has
differential $sI-\bigwedge^rD_j$ on the swept cells.
The fibre-to-cap map is the identity on $\bigwedge L_0$;
on the adapted long-circle factor it is the integral
norm matrix $N_{m_jn_j}(s,D_j)$ determined by
\[
 (sI-D_j)N_{m_jn_j}(s,D_j)=(h-1)I,
 \qquad s=a\text{ or }b.
 \tag{E.C14}
\]
For positive $n_j$ the norm is
$\sum_{r=0}^{m_j-1}s^rD_j^{m_j-1-r}$, and for
$n_j=-1$ it is minus $s^{-m_j}$ times that sum.
In phase degree $r$, use
$N_{m_jn_j}(s,\bigwedge^rD_j)$, rather than the
exterior power of the norm matrix.
The boundary base-circle generator maps to the cap circle, while
the cell containing both that meridian and the adapted
long circle maps to zero. These are actual winding
maps in the mapping-torus coordinates. The linear
three-torus return fixes its origin; hence there is
no additional translation prism in this comparison.

Form the core over a bouquet of three circles with
deck images $a,b,1$ and monodromies $A_1,A_2,M_0$.
Attach the two finite complexes by (E.C14) and the
cusp by (E.C12), using their double mapping cylinders.
Through degree three the resulting free $R$ complex
has ranks $(4,20,49,62)$. Its entries and every unit
cancellation are given exactly in
\texttt{code/maps\_complement\_module.py} and its
JSON output. The program uses the integral normal
form $h^r$ times an alternating word in
$a,a^2$ and $b,b^2,b^3$; no finite group quotient or
floating-point approximation is used.

For completeness the final reduction can be checked
without reading a long matrix. Cancelling degree-zero
and degree-one unit pairs leaves ranks $(1,2,34,62)$,
with $d_1=(a-1,b-1)$. The only nonzero columns of
$d_2$ are $F$ and $F(b-1)$, where
\[
 F=\begin{pmatrix}1+a+a^2\\h(1+b+b^2+b^3)\end{pmatrix}.
 \tag{E.C15}
\]
The first component of $F$ is injective on $R$:
multiplication by $a-1$ converts it to $h-1$, and
$R$ is free over the Laurent ring $A$. Thus removing
the $F$ column after subtracting its multiple gives
a free basis of $\ker d_2$ of rank33. The corresponding
row of $d_3$ is zero by the chain equation. Unit
reductions of the remaining presentation leave five
generators $(x,t,y,z,e)$. One generator $e$ occurs
in no relation. Three pure $t$ relations are
\[
 tq=0,\qquad t(1-a^{-1})=0,
                  \qquad t(b-a^{-1})=0.
 \tag{E.C16}
\]
They generate exactly the augmentation relations on
$t$: multiply the second by $a$ to get $t(a-1)=0$,
and subtract it from the third to get $t(b-1)=0$.
Every other pure $t$ coefficient has augmentation
zero and is consequently redundant. Reducing the
remaining $t$ coefficients by augmentation, three
relations give successively $yq=0,zq=0,xq=0$.
The final three relations are
\[
 x(1+b)+3t=0,\quad x(b+b^2)+3t=0,
                   \quad x(1+b+b^2+b^3)+6t=0.
 \tag{E.C17}
\]
The last two are respectively the first multiplied
by $b$, and the first multiplied by $1+b^2$.
These identities prove (E.C1). The program verifies
the original chain equations, all Schur complements,
and the coefficient list used in (E.C16)--(E.C17).

The invariant subgroup generated by $t$ injects in
$L$. One way to see this is to start with the exact
sequence of right $\mathbb ZC_4$ modules
\[
 0\longrightarrow\mathbb Z\longrightarrow
 \mathbb Z[C_4/C_2]\longrightarrow\mathbb Z_-\longrightarrow0,
\]
where the first map is the sum of the two cosets.
Induce to $Q$ and push out its first term by the
augmentation to the trivial module $\mathbb Z$.
The resulting presentation is (E.C1), after the
change $x\mapsto x+t$ which replaces the coefficient
$3$ by $1$. Pushout preserves the displayed injection.
Restriction to $C_4$ shows nonsplitting: an invariant
correction of the lift $x$ changes this odd coefficient
by an even number. Equivalently the standard cyclic
resolution computes
$\operatorname{Ext}^1_{\mathbb ZC_4}(\mathbb Z_-,\mathbb Z)
=\mathbb Z/2$. Induction--restriction gives the same
extension group for (E.C3).

Finally insert the actual fibre cycles $uw,ud,wd$ in
the core before any cancellation and carry their
coordinate columns through the same basis changes.
The five-generator presentation gives respectively
\[
 t,\qquad t(a^{-1}-a),\qquad t(a^{-1}-1).
 \tag{E.C18}
\]
Relations (E.C16) reduce these to (E.C4). The external
boundary cover has homotopy type $T^3$ and these are
its three degree-two generators. This also checks
the ordinary coefficient independently: after taking
$K$ coinvariants, $2x+3t=0$, so $t$ is twice a
primitive class, exactly the coefficient2 in E.2.6.

To retain the remaining peripheral data, use the
boundary word $abc$, where the cusp generator $c$
has deck image1. On swept fibre chains its map to
the three core edges has blocks
\[
                 bI,\qquad F(A_1),\qquad F(A_1)F(A_2).
 \tag{E.C21}
\]
The identity
\[
 (aI-F(A_1))b+F(A_1)(bI-F(A_2))
       +F(A_1)F(A_2)(I-F(M_0))=(ab-1)I
\]
checks its chain equation. This is the usual prism
subdivision of the actual peripheral word, with
its given ribbon marking. Tracking all ten
degree-two columns through the same reductions,
and using only $t(a-1)=t(b-1)=0$ and
$xq=yq=zq=0$, gives (E.C19a)--(E.C19b).
The unreduced columns are saved in the JSON
certificate. In particular no assertion that a
homologically invisible peripheral component is
zero enters this calculation.
\end{proof}
